\documentclass[11pt]{article}

\usepackage[T1]{fontenc}
\usepackage[utf8]{inputenc}
\usepackage[margin=1in]{geometry}
\usepackage{hyperref}
\hypersetup{colorlinks=true, linkcolor=blue, urlcolor=blue, citecolor=blue}
\usepackage{parskip}
\usepackage{enumitem}

% ---------------------------------------------------------------------
% Formatting helpers
% ---------------------------------------------------------------------
\newcommand{\origcomment}[1]{\textit{``#1''}}
\newcommand{\responsespace}{%
  \par\medskip\noindent\textbf{Response:}\par
  \vspace{3.2cm}%
  \par\noindent\hrulefill\par\bigskip
}
\newcommand{\filledresponse}[1]{%
  \par\medskip\noindent\textbf{Response:} #1\par
  \vspace{0.4cm}%
  \par\noindent\hrulefill\par\bigskip
}
\newlist{commentlist}{itemize}{1}
\setlist[commentlist]{leftmargin=2.2cm, itemsep=1.6em, label=\textbf{[label]}}

\title{Response to Editor and Reviewer Comments}
\author{}
\date{}

\begin{document}
\maketitle
\vspace{-2.5cm}

\begin{center}
\begin{tabular}{ll}
\textbf{Manuscript title:} & HR-VILAGE-3K3M: A Human Respiratory Viral Immunization \\
 & Longitudinal Gene Expression Dataset for Systems Immunity \\[4pt]
\textbf{Journal:} & Scientific Data \\
\textbf{Submission ID:} & acf167d6-60f5-442f-9760-05d400c4dabd \\
\end{tabular}
\end{center}

\bigskip
\noindent We thank the Editor and both reviewers for their careful reading of the manuscript and
for the constructive comments below. We provide a point-by-point response to every comment.
Line/page numbers below refer to the revised manuscript unless otherwise noted.

\bigskip
\noindent \textit{[Optional: insert a short summary of the overall revision here, e.g. that the
manuscript has been restructured as a Data Descriptor, a Data Records section has been added,
the Related Work and Potential Applications sections have been removed/condensed, and references
have been switched to a numbered format.]}

\vspace{1cm}
\hrule
\vspace{1cm}

% =======================================================================
\section*{Response to the Editor}
% =======================================================================

\subsection*{Specific Editor Comments Requiring a Response}

\begin{commentlist}

\item[E1.] \origcomment{Please note, the 2018 ImmPort publication in Scientific Data was not a
data descriptor. We would not be able to consider this work under our article format, only as a
Data Descriptor describing a new, integrated resource that readers would be able to download and
utilize in their own work. The goal is to share data for others to use, with the details in the
paper intended to provide transparency (how data were collected or validated) or practical help
(how to find the data, what variables/column headings mean, etc.). Please do not add any
additional analyses or results that are out of scope of a Data Descriptor.}
\filledresponse{Understood, and agreed. We have revised the manuscript throughout to conform to
the Data Descriptor format, and have removed content that read as general results, analyses, or
conclusions about what the data show, retaining only content needed for transparency (Methods) or
practical use (Data Records, Usage Notes) of the resource. Concretely: the dataset name has been
removed from the title (E4); section numbering has been removed (E2); references now use a
numbered format (E3); the Background and Summary section has been trimmed to motivation and prior
art only, with the former ``Related Work'' section condensed into it and its results-style
sentences removed (E5); a new Data Records section describing the repository's file/folder
structure and column definitions has been added (E10); a Usage Notes section with practical
data-loading examples has been added; and the ``Potential HR-VILAGE-3K3M Applications'' section
has been removed in full (E13). The Technical Validation section is being revised along the same
principle -- QC/validation content only, no proof-of-utility framing -- as detailed in our
response to E12 below.}

\item[E2.] \origcomment{Remove the section numbering.}
\filledresponse{Done. Section numbering has been removed throughout the manuscript
(\texttt{secnumdepth} set so no section, subsection, or subsubsection number prints). All
headings in the main text and Supplementary Information (e.g., ``Background and Summary,''
``Methods,'' ``Technical Validation'') now appear without numeric prefixes.}

\item[E3.] \origcomment{Please use a numbered reference format}
(\url{https://www.nature.com/sdata/publish/submission-guidelines#references}).
\filledresponse{Done. References have been switched from the author-year format to a numbered
format: in-text citations now appear as bracketed numbers in order of first citation (e.g.,
[1], [2]), and the reference list is numbered accordingly in that same citation order.}

\item[E4.] \origcomment{Remove `HR-VILAGE-3K3M' from the manuscript title.}
\filledresponse{Done. The dataset name has been removed from the title. The title now reads:
``A Human Respiratory Viral Immunization Longitudinal Gene Expression Dataset for Systems
Immunity.''}

\item[E5.] \origcomment{Please check the Background and Summary section to ensure it is in scope
and does not contain any general results, analyses, conclusions or facts relating to what the
data shows. The section should be limited to brief descriptions of the motivations of the project
and any relevant prior art. Data Descriptors do not analyse data or present findings, as the goal
is for others to analyse the data for themselves. Please also condense the `related work' section
into the Background and Summary section.}
\filledresponse{Done. The Background and Summary section has been revised to contain only
motivation and prior art. We removed the paragraph describing potential downstream applications
(biomarker discovery, AI/ML model development, etc.), which read as a general conclusion about the
data rather than motivation, and condensed the paragraph describing our curation effort into a
brief summary retaining the key elements of the effort (identifier and gene-symbol harmonization,
obtaining missing antibody titers from study authors, and quality control) while pointing to the
Methods section for full procedural detail. The former ``Related Work'' section (its two
subsections on existing transcriptomic resources and on batch-effect correction/harmonization
methods) has been condensed into a single paragraph and merged into Background and Summary,
retaining the citations that establish prior art and the gap this resource fills (including an
additional cross-platform batch-correction reference, Sun et al., bioRxiv 2026,
doi:10.1101/2024.09.30.615938), while removing
the forward-looking, results-style sentence (``In this work, we benchmark multiple normalization
strategies\ldots''). The remaining motivation and challenges paragraphs were also tightened for
concision. The section is now approximately 580 words, within \textit{Scientific Data}'s 700-word
limit for Background and Summary. A color-coded tracked-changes version of this edit is provided as
\texttt{manuscript/template\_trackchanges.tex} alongside the clean revised manuscript.}

\item[E6.] \origcomment{Because at least some of your data looks to originate from third party
datasets, either shared directly or used as input data to compute or train a new output, this
policy applies, we ask some general questions:}

\textbf{Q1.} \origcomment{Please tell us how you obtained the input data and where from, sharing
enough detail to allow someone to find the same records. This could include sharing any input
data DOIs, URLs, accession numbers, or specific search queries for third party platforms. If data
were not previously available for public download we need to know the manual process for how you
obtained it, such as the name of the organization, contact details, and method of the
collaboration, to allow someone else to obtain this via the same route.}
\filledresponse{Gene expression data were obtained from three publicly available repositories:
GEO, ImmPort, and ArrayExpress. The original collection was assembled via targeted search queries
against these repositories, as detailed in Lan et al.~(2024); HR-VILAGE-3K3M extends this
collection with additional studies identified using the same targeted-search protocol against
records added since that collection. Metadata (outcome variables, subject characteristics) were
obtained from GEO, ImmPort, and the supplementary files of the original publications.
Subject-level antibody titers not already deposited in GEO/ImmPort or included in supplementary
files were requested directly from the corresponding authors of the original publications via
email. Accession numbers, repository names, and direct URLs for all 66 constituent studies are
provided in Supplementary Table 1, allowing every record to be located directly.}

\textbf{Q2.} \origcomment{Please also confirm that information from Q1 is included in the Methods
sections (e.g. state the line number).}
\filledresponse{Yes. This information is stated in Methods, ``Data collection and curation''
paragraph (\texttt{template.tex}, line 252), which describes the repositories and search protocol
used to assemble the gene expression and metadata, the process for obtaining subject-level
antibody titers, and now also references Supplementary Table 1 for the complete list of accession
numbers and repository links for all 66 studies.}

\textbf{Q3.} \origcomment{Please tell us how you ensured you could re-use the data. For data from
private sources, please confirm the data owner has provided permission for you to re-share these
data under CC0 or CC-BY licenses (or data usage agreement for identifiable human data). For data
from public (open / downloadable) sources, please share a link to the license or terms that
confirm this. Or, if data are not within the scope of copyright at all, confirm this applies
and/or your re-distribution is acceptable use. Finally, if you are a tertiary user of a secondary
dataset that is itself sourced from elsewhere, we still request some basic validation is
performed to ensure the primary source was open and legitimately re-distributed.}
\filledresponse{For data obtained from GEO, ArrayExpress, and ImmPort, all three repositories
confirm re-use rights for depositor-submitted data: NCBI states that it places no restrictions on
the use or distribution of GEO data
(\url{https://www.ncbi.nlm.nih.gov/geo/info/disclaimer.html}); EMBL-EBI's Terms of Use impose no
additional restriction on ArrayExpress data beyond those set by the data owner
(\url{https://www.ebi.ac.uk/about/terms-of-use}); and the ImmPort Data Use Agreement grants the
unrestricted right to use and distribute ImmPort data, provided redistribution occurs under
commensurate terms (\url{https://docs.immport.org/home/agreement}). \textbf{[TO DISCUSS: re-use
confirmation for subject-level antibody titers obtained directly from study authors via
email --- to be added once resolved.]}}

\item[E7.] \origcomment{Is the data at Lan et al. (2024) peer-reviewed and published? If not,
please incorporate all data into this work formally. Otherwise, please ensure it is transparent
in the manuscript text (and response letter) what specific new data has been added to this work
over that of previous publications. This will be assessed against our related publication policy}
(\url{https://www.nature.com/sdata/policies/editorial-and-publishing-policies#complementary}).
\filledresponse{Lan et al. is now peer-reviewed and published: F1000Research 2025, 14:493
(\url{https://doi.org/10.12688/f1000research.162267.1}), first published 13 May 2025, status
``version 1; peer review: 1 approved with reservations,'' distributed under a CC-BY license. The
bibliography entry has been updated accordingly (from the ``awaiting peer review'' status cited in
our earlier preprint to the current published, peer-reviewed status).

Lan et al. is a dataset-discovery and curation paper: it systematically searched GEO, ImmPort, and
ArrayExpress and compiled a list of 55 relevant datasets (18 inoculation, 37 vaccination;
approximately 2,513 subjects total), but does not reprocess, harmonize, or host the underlying
expression data itself --- it identifies and describes existing third-party datasets for others to
download and analyze independently. HR-VILAGE-3K3M is a distinct data product built on top of
that search process: it extends the collection to 66 studies and over 3,000 subjects (additional
recent studies, broader ``mixed exposure'' scope beyond the inoculation/vaccination dichotomy, and
curated subject-level antibody measurements obtained directly from study authors), and, critically,
delivers a harmonized, quality-controlled, ready-to-use gene expression and metadata product (with
standardized gene symbols, aligned sample identifiers, and standardized outcome definitions across
studies) rather than a list of source datasets. This distinction, and the specific numeric
comparison to Lan et al., is now stated explicitly in Methods, ``Data collection and curation''
paragraph (\texttt{template.tex}, line 252).}

\item[E8.] \origcomment{Please see our section for citing data -- given only 66 studies are
utilized as input, please use the formats outlined at the following link to cite the input data
at source} (\url{https://www.nature.com/sdata/publish/submission-guidelines#dataoverview}).
\origcomment{Please also ensure to cite the paper(s) that these data came from.}
For GEO data: \textit{Lastname, A.B, Lastname, C.D. Lastname, E.F. GEO.}
\url{https://identifiers.org/geo/GSEXXXXXX} \textit{(20XX)}. For ArrayExpress:
\textit{Lastname, A.B, Lastname, C.D. Lastname, E.F. ArrayExpress.}
\url{https://identifiers.org/arrayexpress:E-MTAB-XXXXX} \textit{(20XX)}.
\filledresponse{Both citation requirements are now addressed for the 59 constituent microarray and
RNA-seq studies (the remaining 7 scRNA-seq studies will be completed in the same way). (i)
\emph{Paper citations:} rather than being left only in a supplementary list, the source
publication(s) for each of these studies are now cited in-text at their point of use, in Methods,
``Data collection and curation'' paragraph (\texttt{template.tex}, line 252), and included in the
numbered reference list via \texttt{bibliography.bib}. (ii) \emph{Data citations:} formal,
Nature-format data citations for each of the 66 constituent studies --- \textit{Lastname, A.B. et
al. GEO.} \url{https://identifiers.org/geo/GSEXXXXXX} \textit{(20XX)} for GEO studies and the
corresponding ArrayExpress/ImmPort forms otherwise --- are provided in the \texttt{data\_citation}
column of Supplementary Table 1, generated from each study's own GEO/ArrayExpress/ImmPort
contributor and submission metadata rather than invented placeholders.}

\item[E9.] \origcomment{Please do not replicate the original input data files at Huggingface, only
the data product that has been developed for this work and any new data that was not previously
available should be present at the repository. Please simply cite the inputs so that readers will
still go to the original resources to obtain these data.}
\filledresponse{Done. The raw (unprocessed) input files are being removed from the Hugging Face
repository; only the harmonized, quality-controlled data product (processed expression matrices,
standardized metadata, and curated antibody measurements) is retained there, together with any
newly created data (e.g., curated subject-level antibody titers obtained directly from study
authors) that was not previously available elsewhere. Readers seeking the original raw files are
directed to the source repository for each study rather than to a local copy: the
\texttt{study\_meta.csv} index and Supplementary Table 1 cite each of the 66 constituent studies at
its original GEO, ArrayExpress, or ImmPort accession, from which the raw files can be obtained
directly. This is now stated in the Data Records section (\texttt{template.tex}, opening
paragraph). \textbf{(Ruilie, please delete the raw data from the Hugging Face repository and add
the study-level supplementary table to Hugging Face.)}}

\item[E10.] \origcomment{The Data Records section is missing. Please add a Data Records section
and ensure it provides a description of the files being shared. This should start with an opening
sentence such as ``The dataset is available at HuggingFace (reference number)'' followed by a
practical description of the files, including file names and folder structure (where these are
not obvious or simple) and a description of column headings or variable names (where not
implicit). No other content, such as a statistical summary of the data, or further methods (place
these in the methods section), or detailed analyses should be included. If you wish to include
any summary statistics, please do so in a ``Data overview'' section.}
\filledresponse{Done. A new Data Records section has been added between Methods and Technical
Validation, opening with a sentence naming the Hugging Face repository
(\texttt{xuejun72/HR-VILAGE-3K3M}), its URL, and its DOI-based data citation (see E11 below). The
section then describes the repository's file and folder structure --- the top-level
\texttt{study\_meta.csv} index, the four per-study data folders (\texttt{bulk\_gene\_expr/},
\texttt{single\_cell\_gene\_expr/}, \texttt{meta/}, \texttt{antibody/}; raw input files are not
mirrored in the repository, per E9 below), and the supporting reference files
(\texttt{Meta\_variable\_description.csv},
\texttt{example\_dataset/}, \texttt{croissant.json}, \texttt{relation.png}) --- along with the
column headings for \texttt{study\_meta.csv} and the \texttt{meta/} files. Consistent with this
comment, no statistical summaries or additional methods detail were included in this section; the
existing statistical summary at the end of Methods (participant/study counts, platform and
sample-type breakdowns) was left in place there rather than duplicated or moved, since it is
methods-adjacent descriptive summary rather than a new result.}

\item[E11.] \origcomment{Please add a data citation for the dataset to the reference list using
these instructions} (\url{https://www.nature.com/sdata/submission-guidelines#data_citations}
-- \origcomment{note that a DOI URL should be used, not any other link}).
\origcomment{Please add the reference number to wherever the dataset is mentioned in the text --
the main position should be the first part of the Data Record in a sentence describing where the
data has been deposited.}
\filledresponse{Done. Hugging Face has assigned the dataset a DOI (\url{https://doi.org/10.57967/hf/9892}),
minted via the Hub's built-in DOI/DataCite integration. A corresponding data citation has been
added to the reference list (\texttt{bibliography.bib}), and is referenced by number at the
opening sentence of the Data Records section (\texttt{template.tex}, line 288), which now reads
``...deposited in the Hugging Face Datasets repository \texttt{xuejun72/HR-VILAGE-3K3M}...(data
citation: [reference number]).'' The DOI is also given in the Data Availability statement.}

\item[E12.] \origcomment{Please check and revise the current Technical Validation section to make
sure it only contains a summary of how the data were checked or validated. This could include
statistical checks, comparisons with other data, further discussions of methods or limitations,
or any other process for quality control. This section should not include general results,
discussion or analyses -- data analyses are out of scope for Data Descriptors unless they are
relevant for checking its quality or integrity. Proofs that data can be used for specific use
cases are not required -- because it is not a requirement that Data Descriptors prove utility for
specific use cases -- if the data can be considered generally sound, so please consider general
checks that will be helpful for most users.}
\filledresponse{Agreed, and revised. We reviewed each Technical Validation subsection against this
standard, benchmarking our approach against comparable multi-study transcriptomic Data Descriptors
(e.g., the Immune Signatures Data Resource, \textit{Sci. Data} 2022), which similarly limit
Technical Validation to checks on the data's technical soundness (QC metrics, batch-effect
assessment, metadata-imputation accuracy) with no predictive-modeling or ``which-method-wins''
benchmarking. Age and Sex Prediction is retained but reframed as a QC screen only (accuracy of
imputed sex/age against the recorded metadata, used to flag potential sample mislabeling), removing
the secondary framing of it as a cross-platform predictive-modeling benchmark. The subsection
combining batch-effect correction with an antibody-responder prediction benchmark (RNN/LSTM/GRU/
Transformer models, AUC/accuracy/F1 comparisons) is being cut down to its batch-effect visualization
check only -- i.e., confirming that harmonization brings samples from different studies into
alignment while preserving biological (responder-status) separation -- since the downstream
predictive-modeling comparison is exactly the ``proof of utility for a specific use case'' this
comment says is not required, and is also the subject of Reviewer 2's comment 5 regarding possible
QN artifacts in the reported AUCs. Cell Type Annotation using paired bulk and single-cell data is
retained, as it validates internal consistency between two modalities measured on the same
subjects rather than demonstrating a downstream application. \textbf{[PENDING: the batch-effect-only
reanalysis is being rerun to regenerate the corresponding figure/results before this is reflected
in \texttt{template.tex}.]}}

\item[E13.] \origcomment{Remove the section ``5 Potential HR-VILAGE-3K3M Applications''.}
\filledresponse{Done. The section (change point detection, gene expression imputation, causal
inference, deconvolution, building prediction models, and pretraining/fine-tuning) and its
accompanying figure have been removed in full.}

\item[E14.] \origcomment{HuggingFace is not set up as a repository we generally accept. Please
adjust the license to CC-BY, and set up a DOI for the repository. In addition, please ensure the
metadata for the repository is completed -- generally this will involve setting up authors,
title, etc. Please then also update the dataset card to note the citation for the dataset so that
this is easier for readers to find.}
\filledresponse{A DOI has been generated for the repository via Hugging Face's built-in
DOI/DataCite integration (\url{https://doi.org/10.57967/hf/9892}; see E11). \textbf{(Ruilie,
please modify the license to CC-BY, complete the repository metadata (full author list, title),
and update the dataset card to display the DOI-based citation.)}}

\end{commentlist}

\subsection*{General Comments}
\noindent\textit{These are general notes communicated for all manuscripts and, per the Editor, do
not require an individual written response. Listed here for tracking/checklist purposes only.}

\begin{itemize}[itemsep=0.3em]
  \item One clean (unhighlighted, no tracked changes) copy of the main article file in a
  machine-readable format (\texttt{.docx} or \texttt{.tex} only, no PDF;
  \url{https://www.nature.com/sdata/submission-guidelines#latex}).
  \item [Optional] A tracked-changes/highlighted version of the paper as a related manuscript
  file.
  \item All figure images, one file per figure (merged panels where applicable).
  \item Supplementary information as a PDF only, if not merged into the main paper.
  \item Any oversize tables (\textgreater\,1 A4 page) as Supplementary Tables in \texttt{.xlsx} or
  \texttt{.csv} format.
  \item This single ``response to comments'' file, answering all Editor and Reviewer comments.
  \item Confirmation that all datasets are live (not under embargo) and formally cited
  (\url{https://www.nature.com/sdata/policies/data-policies#publication}).
\end{itemize}

\clearpage
% =======================================================================
\section*{Response to Reviewer 1}
% =======================================================================

\noindent\textit{[Reviewer 1 summary: the authors have developed a metadata resource integrating
bulk and single-cell transcriptomic profiles from 3,178 subjects across 66 studies; the reviewer
notes the resource is not comprehensive and appears to have missed several key datasets, and asks
that the key issues below be addressed prior to publication.]}

\bigskip
\subsection*{Major Comments}

\begin{commentlist}
\setlist[commentlist]{leftmargin=2.2cm, itemsep=1.6em, label=\textbf{[label]}}

\item[R1-1.] \origcomment{Authors should provide PRISMA workflow describing data collection
process.}
\responsespace

\item[R1-2.] \origcomment{Most of these vaccine studies are done on adults \textless\,60 years.
Authors should explicitly show how many samples are coming from the older adults (Age
distribution of the sample for longitudinal time points) -- for both scRNAseq, bulk RNAseq and
microarray.}
\responsespace

\item[R1-3.] \origcomment{What is sample ratio based on vaccine type (Flu IAV, LIAV, Fluzone HS,
FluAd, Flublok, FLuarix, similarly COVID vaccines and others)?}
\responsespace

\item[R1-4.] \origcomment{Define Inoculation and mixed exposure.}
\responsespace

\item[R1-5.] \origcomment{How are single cell data dealt? How are cell types defined -- clarify on
those details.}
\responsespace

\item[R1-6.] \origcomment{Donors with high baseline titers mask seroconversion rates. How are
these donors dealt with? Authors should consider using baseline adjusted methods for response
categorization (maxRBA from Steve Kleinstein's group).}
\responsespace

\item[R1-7.] \origcomment{How are they defining response status for other vaccines?}
\responsespace

\item[R1-8.] \origcomment{Authors should validate the overall processing data by checking known
vaccination read outs -- interferon signature at day 1, plasmablast expansion at day 7 for
influenza vaccine. Authors should consider checking these signatures at these time points in case
of transcriptome cell types.}
\responsespace

\end{commentlist}

\clearpage
% =======================================================================
\section*{Response to Reviewer 2}
% =======================================================================

\noindent\textit{[Reviewer 2 summary: the paper presents a comprehensive respiratory viral
transcriptomics dataset integrating microarray, bulk RNAseq, and scRNAseq data from 66 independent
studies and 3000+ samples. Strengths noted: curation/harmonization effort across 66 studies
(identifier resolution, antibody data outreach, duplicate removal); public availability via
Hugging Face; demonstrated utility via prediction and cell-type annotation applications.]}

\bigskip
\subsection*{Major Comments}

\begin{commentlist}
\setlist[commentlist]{leftmargin=2.2cm, itemsep=1.6em, label=\textbf{[label]}}

\item[R2-1.] \origcomment{(compare with ImmPort) Explicit positioning of this resource amidst
related resources like ImmPort or HIPC (and associated visualizer ImmuneSpace) should be
mentioned in Introduction. Without this explicit positioning, a reader would find it hard to
assess the extra value added by this resource atop existing resources like ImmPort. For instance,
ImmPort had compiled 300 studies based on its 2018 publication in Scientific Data. The current
work takes 66 studies, probably many of them from ImmPort, relevant to respiratory viral
immunization (including both influenza and COVID) and harmonizes them. Within this context,
answering the questions below directly in the paper can clarify the value addition of the
HR-VILAGE resource: (i) What fraction of studies in HR-VILAGE are taken from ImmPort? (ii) Of the
studies taken from ImmPort, which harmonization or data cleaning was missing in ImmPort, and what
curation/analysis has been additionally done in HR-VILAGE-3K3M to address these harmonization
issues? For instance, it is hard to believe that ImmPort did not already handle gene ID issues
like SEPTIN14. (iii) For the studies in HR-VILAGE taken from GEO and ArrayExpress, was more effort
required to perform data cleaning and harmonization (compared to studies taken from ImmPort)?}
\responsespace

\item[R2-2a.] \origcomment{(other modalities) A key limitation of HR-VILAGE-3K3M compared to
ImmPort is its restriction to transcriptomic/antibody-response data, versus ImmPort's additional
support for flow cytometry, cytokine, CyTOF, and related datasets. What is the rationale for this
restriction, given that systems immunology research has largely advanced in recent times due to
integrative analysis of multi-modal data collected from immune-related human cohorts? Won't this
restriction limit the utility of HR-VILAGE for future systems immunology analyses or
re-analyses? This limitation either needs to be acknowledged, alongside the advantages of
HR-VILAGE over ImmPort; or addressed by including analysis-ready data pertaining to some of these
extra modalities in HR-VILAGE.}
\responsespace

\item[R2-2b.] \origcomment{A related point on missing modalities is different assays of
immune/antibody response such as MN (microneutralization) titers, HAI, and ELISPOT. There is
mention of HAI in the paper, but the paper is not clear about whether other assays such as MN and
ELISPOT are stored in HR-VILAGE. Are there studies that have jointly assayed these different
immune endpoints? If so, it would be useful to store these different endpoints in HR-VILAGE to
enable alternative definitions of antibody-responders (continuous-valued as well as binary) to
compare/contrast these outcome definitions from biological and statistical viewpoints.}
\responsespace

\item[R2-3.] \origcomment{(compare with other repo) Comparative analysis of this resource with
other independent state-of-the-art resources is missing. While all studies in HR-VILAGE may not
be available in other independent resources, certain studies may also be present in other
transcriptomic resources such as recount2/3 and Gene Expression Commons (GEXC). For such common
studies, what is the qualitative difference between the data preprocessing pipelines employed by
the different resources, and can the authors quantify the similarities and differences in these
resources' processed datasets and downstream results?}
\responsespace

\item[R2-4a.] \origcomment{(why QN, explain AUC) What was the rationale for selecting quantile
normalization (QN) over alternative approaches for RNAseq harmonization? Was QN performed in a
class-specific manner (i.e., stratified by responder group as high, medium, low) rather than
across the pooled dataset? Prior work has shown that whole-dataset QN can distort true biological
differences, particularly when class-effect proportions are high, whereas class-specific QN
better preserves such differences (Zhao, Wong \& Goh, Sci Rep, 2020:
\url{https://www.nature.com/articles/s41598-020-72664-6}).}
\responsespace

\item[R2-4b.] \origcomment{On a related point, could the authors clarify whether the Transformer
model used for antibody responder prediction was trained directly on QN-normalized gene
expression values? If so, how much of the Transformer's predictive performance (AUC = 0.737) can
be attributed to genuine biological signal vs. artifacts of the QN-induced rank-based rescaling?}
\responsespace

\item[R2-5a.] \origcomment{(study specific pipeline) While standardized/uniform (pre)processing
pipelines employed by HR-VILAGE is preferable to reduce methodological variations, another viable
strategy is to preprocess datasets from different studies in a highly study-specific manner to
account for study-specific artifacts. Given this context, it is highly advisable for HR-VILAGE to
not only provide the standardized analysis-ready data for a study, but also the analysis-ready
data provided by the study's authors whenever available (e.g., from websites/GitHub linked to
papers such as illustrated here: \url{https://www.tsanglab.org/resources}).}
\responsespace

\item[R2-5b.] \origcomment{Conversely, focusing on RNAseq data, gene counts data provided by study
authors is used for certain studies in HR-VILAGE, and gene counts from standardized analysis
(alignment of fastq files and counts summarization) is used for other studies. For each study,
both study-author-provided and uniformly-processed pipeline based versions of the gene counts
matrix should be provided -- i.e., even when RNAseq-based gene counts are available for certain
studies, HR-VILAGE should also provide gene counts recomputed from its standard pipeline. The
outcome would be two versions of the resource (study-specific processing and uniform processing),
useful for QC comparison (e.g., tSNE/PCA) to flag samples with discrepancies between the two
pipelines.}
\responsespace

\item[R2-6.] \origcomment{(compare application results with publications) For any of the three
applications shown, please compare your results with previously published results to indicate
what is common, and more importantly, what is different, so that there is motivation to do
reanalysis of these studies, or reinterpret one study in the context of others (as in a
meta-analysis) to get novel insights. Presenting conserved vs. novel findings, e.g. via simple
Venn diagram overlaps of marker genes predictive of certain outcomes, may suffice to encourage
researchers to use this resource for reanalyzing existing viral transcriptomic datasets.}
\responsespace

\end{commentlist}

\subsection*{Minor Comments and Typos}

\begin{commentlist}
\setlist[commentlist]{leftmargin=2.2cm, itemsep=1.6em, label=\textbf{[label]}}

\item[R2-m1.] \origcomment{Please cite references for standard scRNAseq processing methods/tools,
and expand certain symbols like log1p (transformation).}
\responsespace

\item[R2-m2.] \origcomment{Why repeat the description of dataset preprocessing, transformation,
and related steps, given originally in Methods, again one more time in Results (e.g., Results
section text on sex prediction talks about QN and log1p, etc.)? Could you not say in the Results
section that these tasks use analysis-ready data, preprocessed/normalized as explained in the
Methods section?}
\responsespace

\item[R2-m3.] \origcomment{This sentence seems incomplete: ``with AUCs ranging from 0.602 to
0.702, while ComBat and Regression corrections generally produced lower scores.'' Please fix it.}
\responsespace

\item[R2-m4.] \origcomment{Figure 2 caption says ``HV-RIGEL-3K'' Dataset instead of
HR-VILAGE-3K3M. Please fix it.}
\responsespace

\item[R2-m5.] \origcomment{Day 0 is just before vaccination, day -1, -7 are days prior to
vaccination (or is it pre1 and pre7), etc. -- are these terms standardized across studies?}
\responsespace

\item[R2-m6.] \origcomment{Of the samples drawn from the 66 original studies, were all available
samples retained in the final harmonized dataset, or were some samples excluded based on
quality-control criteria (e.g., low read depth, RIN score, etc.)? Please clarify.}
\responsespace

\item[R2-m7.] \origcomment{Text on Page 18 mentions ``six random seeds'' but Figure 2b has no
error bars for the accuracy and F1 scores. Please fix.}
\responsespace

\item[R2-m8.] \origcomment{Foundation models may need the scale of HIPC ImmPort (300+ studies).
With few studies in HR-VILAGE, can foundation models be really trained?}
\responsespace

\item[R2-m9.] \origcomment{In Section 4.3, what is the overlap of the number of subjects on Day 1
and Day 7 from the datasets GSE201533 and GSE201534? As per metadata GSE201534 has only 6
subjects.}
\responsespace

\item[R2-m10.] \origcomment{Instead of running QC to flag problematic samples, can you also fix
sample mislabeling or other such issues using approaches such as MODMatcher (PLOS Comput Biol,
Yoo et al. 2014)? If not fixing these issues, can some diagnostics at least be done to see if
problematic samples are scattered across studies, or coming exclusively from one or two studies?}
\responsespace

\item[R2-m11.] \origcomment{Methodological details missing: (i) Could the authors specify the
reference genome build (e.g., hg38) and the corresponding Gencode annotation version used for
transcript quantification/alignment? (ii) What guidelines are used to select the datasets to
include (e.g., search terms in PubMed, inclusion/exclusion criteria of papers, etc.)? Targeted
search queries in Lan et al. 2024 are not provided here, and the link to the Lan et al. 2024
F1000 paper is also broken. (iii) For outcome prediction, please add random/naive classifiers'
F1 score or AUC to all plots, so that a baseline is there to interpret other methods' results
(e.g., a classifier that randomly predicts the class proportional to the fraction of positive
samples within each study). Optionally, SHAP or other analysis to see what features are
contributing to prediction accuracy could also help distinguish biological signals vs. other
artifacts.}
\responsespace

\end{commentlist}

\end{document}
